\section{Introduction}

Character performance is shaped by more than the words a character speaks. In conversational animation, an animator may consider who is speaking, who is being addressed, what has occurred earlier in the scene, the relationship between characters, and what a character may be attempting to communicate or conceal. These contextual interpretations influence performance choices such as when a character establishes or avoids eye contact, how strongly the head participates in a gaze shift, whether an expression is sustained or suppressed, and how a character reacts while listening. A short utterance such as \emph{“I’m fine,”} for example, may support substantially different performances depending on whether it follows a casual question, an accusation, or an emotionally significant event. Authoring such performances therefore involves not only producing facial motion, but interpreting narrative and interpersonal context and translating that interpretation into temporally coordinated animation decisions.

Procedural and data-driven animation systems have substantially reduced the effort required to produce lower-level character motion. Systems such as JALI generate expressive, editable speech animation from audio and transcripts, while subsequent approaches such as VisemeNet produce animator-centric speech motion curves directly from audio \cite{edwards2016jali,zhou2018visemenet}. More recently, large language models (LLMs) have enabled natural-language interfaces for animation generation and control. Huang et al. demonstrated that LLM-generated structured commands can drive rigged characters \cite{huang2023realtime}; Keyframer combines natural-language generation with animation variants and direct editing \cite{tseng2025keyframer}; and LogoMotion provides code-connected editing widgets for modifying generated motion, grouping, and timing \cite{liu2025logomotion}. These systems demonstrate that LLMs can lower barriers to animation generation while still supporting iterative human refinement.

Recent work has also begun to connect textual context to character facial behavior. Prompt-to-Animation, for example, uses LLMs to map narrative scenarios grounded in the OCC emotion model to FACS Action Units and demonstrates that the resulting expressions can communicate contextual affect without custom model training \cite{durupinar2025prompt}. This suggests that contemporary LLMs can provide a useful source of high-level interpretation in addition to low-level animation control. However, generating an appropriate facial expression from narrative context does not by itself address how an animator authors a performance when multiple interpretations are plausible.

The limitation we address is therefore not that existing AI animation systems are uneditable. Both Keyframer and LogoMotion show the value of combining generation with direct manipulation and localized editing \cite{tseng2025keyframer,liu2025logomotion}. Rather, the object being edited is typically the generated animation, code, or motion parameters. For character acting, an undesirable motion may instead originate from a higher-level disagreement: the system may have inferred the wrong communicative intent, selected the wrong social target, interpreted avoidance as engagement, or assigned an inappropriate relationship between expression and gaze. When this interpretive layer remains implicit, an animator must infer why the system produced a behavior before deciding how to revise it. Editing a gaze trajectory, for example, is different from communicating that the character should avoid eye contact because they are concealing information.

Work on human-AI creative tools provides evidence that exposing intermediate abstractions can support more meaningful intervention. DreamGarden externalizes an LLM-generated hierarchical plan and allows users to modify an otherwise semi-autonomous generation process through operations such as pruning, extending, and providing feedback \cite{earle2025dreamgarden}. Piet similarly demonstrates the value of designing domain-specific interactive representations around the progressive workflows of creative professionals \cite{shi2024piet}. These systems motivate our central question: \textbf{what should an intermediate representation look like when the object being authored is not an implementation plan or a visual parameter, but a character's interpretation and performance within a conversation?}

We explore this question through a context-aware LLM-assisted authoring workflow for conversational character animation. Our approach introduces an intermediate \textbf{Performance Plan} between narrative interpretation and animation execution. We focus on authoring the performance of a target character in a dyadic conversation, while using information from both participants in the interaction. The system considers structured context including dialogue turns, speaker and addressee information, preceding conversational context, character relationships, and the current utterance. From this context, the LLM proposes semantic performance events describing dimensions such as communicative intent, affect, gaze target and timing, gaze avoidance and re-engagement, head involvement, and related facial-performance parameters. Each event can additionally contain a concise, user-facing rationale linked to observable dialogue or scene evidence.

Crucially, the Performance Plan is not treated as a hidden intermediate step on the way to a generated animation. It is presented to the animator as an editable authoring representation. Our interface allows animators to inspect performance decisions alongside dialogue and timeline context, modify individual semantic parameters, preserve or lock satisfactory decisions, and selectively regenerate others. The resulting plan is then translated into procedural facial-animation controls and executed through JALI within Maya. In this workflow, the LLM does not determine a final performance autonomously. Instead, it proposes an interpretation that can be accepted, contested, or revised before being committed to low-level animation.

We investigate this approach through three research questions:

\textbf{RQ1:} How does conversational and narrative context affect the appropriateness of LLM-proposed character performance decisions?

\textbf{RQ2:} How does an editable semantic Performance Plan affect animators' perceived control and their ability to iteratively refine AI-assisted character animation?

\textbf{RQ3:} How do animators accept, modify, reject, and reinterpret AI-generated performance suggestions during the authoring process?

To investigate these questions, we develop a prototype integrating context-aware performance interpretation, an editable Performance Plan, and procedural character-animation generation. We evaluate the system in a within-subject study with \textbf{[N] participants with character-animation experience}, comparing \textbf{[baseline condition]} with the proposed Performance Plan workflow across \textbf{[N] conversational scenes}. We examine authoring behavior including completion time, localized edits, regeneration patterns, and the acceptance, modification, or rejection of AI suggestions, together with perceived control, workload, and qualitative accounts of how participants incorporated or contested AI interpretations. \textbf{[Replace with final study design and headline results after data collection.]}

This work makes three contributions. First, we introduce a domain-specific intermediate representation that externalizes LLM interpretations of conversational character performance as inspectable and editable semantic decisions. Second, we present a mixed-initiative authoring workflow that connects narrative context to procedural facial-animation controls while supporting localized modification, preservation, and regeneration of performance choices. Third, through an empirical study with animators, we characterize how human creators accept, revise, and contest AI-generated performance interpretations, and derive implications for designing human-AI creative tools in domains where generation depends on subjective interpretation rather than a single correct output.
